\EMkeepfig{m02-mapping}{3}

\EMsection{توابع مختلط}{توابع مختلط}

\EMfigure{m02-mapping}{تابع مختلط \EMim{w=f(z)} نقاط دامنه در صفحهٔ \EMim{z} را به نقاط برد در
صفحهٔ \EMim{w} می‌نگارد}

\EMdisp{z\EMbr =x+iy,\qquad\EMbr  w\EMbr =u+iv}
\EMdisp{f:\mathbb{C}\to\mathbb{C},\qquad\EMbr  f:\begin{cases}u=u(x,y)\\ v=v(x,y)\end{cases}}

\begin{EMexample}{}

تابع \EMim{w=z^2} را در نظر بگیرید:
\EMdisp{u+iv\EMbr =(x+iy)^2\EMbr =(x^2-y^2)+i(2xy)\quad\EMbr \EMbr \Longrightarrow\quad\EMbr \begin{cases}u=x^2-y^2\\ v=2xy\end{cases}}

\end{EMexample}

\EMkeepfig{m02-loci}{3}

\EMsection{چند مکان هندسی مهم}{چند مکان هندسی مهم}

\EMfigure{m02-loci}{مکان‌های هندسی مهم: دایره، \EMim{\rho}-همسایگی نقطهٔ \EMim{a} و ناحیهٔ
حلقوی}

\EMdisp{|z|\EMbr =1,\qquad\EMbr  |z-a|\EMbr =\rho,\qquad\EMbr  |z-a|<\rho,\qquad\EMbr  \rho_1<|z-a|<\rho_2}

مجموعهٔ \EMim{|z-a|<\rho} را \EMim{\rho}-همسایگی \EMim{a} می‌نامند.

\EMhold

\EMsection{حد توابع مختلط}{حد توابع مختلط}

\begin{EMdefinition}{حد}

اگر به ازاء هر \EMim{\varepsilon>0} بتوان \EMim{\delta>0} به‌دست آورد، به طوری که از نامساوی \EMim{|z-z_0|<\delta} نتیجه شود \EMim{|f(z)-l|<\varepsilon}، آنگاه
\EMdisp{\lim_{z\to z_0}f(z)\EMbr =l}
بنابراین به ازاء کلیهٔ نقاط داخل دایره‌ای به مرکز \EMim{z_0} و به شعاع \EMim{\delta} لازم است که \EMim{f(z)} داخل دایره‌ای به مرکز \EMim{l} و به شعاع \EMim{\varepsilon} باشد. در اینجا \EMim{\varepsilon} داده می‌شود و \EMim{\delta(\varepsilon)} را بایستی به‌دست آورد. واضح است مفهوم حد در اینجا گسترده‌تر از حد در توابع حقیقی است.

\end{EMdefinition}

\EMfigure{m02-limit}{تعریف حد: دایرهٔ \EMim{\delta}-همسایگی \EMim{z_0} در صفحهٔ \EMim{z} و
دایرهٔ \EMim{\varepsilon}-همسایگی \EMim{l} در صفحهٔ \EMim{w}}

\EMhold

\EMsection{پیوستگی توابع مختلط}{پیوستگی توابع مختلط}

\begin{EMdefinition}{پیوستگی در یک نقطه}

تابع \EMim{w=f(z)} در نقطهٔ \EMim{z=z_0} پیوسته است هرگاه سه شرط زیر برقرار باشد:

\begin{enumerate}
\def\labelenumi{\arabic{enumi}.}
\tightlist
\item
  \EMim{f(z)} در نقطهٔ \EMim{z=z_0} معین باشد.
\item
  \EMim{\displaystyle\lim_{z\to z_0}f(z)} وجود داشته باشد.
\item
  \EMim{\displaystyle\lim_{z\to z_0}f(z)=f(z_0)} باشد.
\end{enumerate}

\end{EMdefinition}

\begin{EMexample}{}

آیا تابع \EMim{f(z)=\sin(1/z)} در نقطهٔ \EMim{z=0} پیوسته است؟

\end{EMexample}

\begin{EMexample}{}

آیا تابع
\EMdisp{f(z)\EMbr =\begin{cases}\sin(1/z)&z\neq0\\ 5+4i&z=0\end{cases}}
در نقطهٔ \EMim{z=0} پیوسته است؟

\end{EMexample}

\begin{EMexample}{}

آیا تابع
\EMdisp{f(z)\EMbr =\begin{cases}3z^3+4&z\neq0\\ 17&z=0\end{cases}}
در نقطهٔ \EMim{z=0} پیوسته است؟

\end{EMexample}

\EMhold

\EMsubsection{خواص توابع مختلط پیوسته}{خواص توابع مختلط پیوسته}

\begin{EMtheorem}{قضایا}

\begin{enumerate}
\def\labelenumi{\arabic{enumi}.}
\tightlist
\item
  اگر توابع \EMim{f(z),g(z)} در نقطهٔ \EMim{z_0} پیوسته باشند، آنگاه \EMim{f(z)\pm g(z)} در \EMim{z_0} پیوسته است.
\item
  اگر توابع \EMim{f(z),g(z)} در نقطهٔ \EMim{z_0} پیوسته باشند، آنگاه \EMim{f(z)\times g(z)} در \EMim{z_0} پیوسته است.
\item
  اگر توابع \EMim{f(z),g(z)} در نقطهٔ \EMim{z_0} پیوسته باشند و \EMim{g(z_0)\neq0} آنگاه \EMim{\dfrac{f(z)}{g(z)}} در \EMim{z_0} پیوسته است.
\end{enumerate}

\end{EMtheorem}

\begin{EMdefinition}{پیوستگی در یک میدان}

تابع \EMim{w=f(z)} در میدان \EMim{D} پیوسته است هرگاه تابع \EMim{f} در تمام نقاط میدان \EMim{D} پیوسته باشد.

میدان: تمام نقاط داخل یک منحنی بسته.

\end{EMdefinition}

\EMfigure{m02-domain}{میدان \EMim{D} در صفحهٔ مختلط}

\EMhold

\EMsection{مشتق توابع مختلط}{مشتق توابع مختلط}

\begin{EMdefinition}{مشتق}

تابع \EMim{w=f(z)} در نقطهٔ \EMim{z=z_0} مشتق‌پذیر است هرگاه حد زیر موجود باشد:
\EMdisp{f'(z_0)\EMbr =\lim_{|\Delta z|\to0}\frac{f(z_0+\Delta z)-f(z_0)}{\Delta z}}

\end{EMdefinition}

\begin{EMremark}{}

وجود مشتق در توابع مختلط نیازمند شرایط سخت‌تری نسبت به توابع حقیقی است، زیرا اگر در صفحهٔ \EMim{z} از \emph{تمام} جهات به نقطهٔ \EMim{z_0} نزدیک شویم، بایستی این حد موجود بوده به مسیر انتخابی بستگی به سمت \EMim{z_0} بستگی نداشته باشد.

\end{EMremark}

\EMfigure{m02-approach}{نزدیک شدن به نقطهٔ \EMim{z_0} از تمام جهات}

\EMhold

\EMsection{شرایط کشی--ریمان}{شرایط کشی–ریمان}

\begin{EMtheorem}{شرط لازم مشتق‌پذیری}

شرط لازم برای آنکه تابع مختلط \EMim{w=f(z)} مشتق‌پذیر باشد آن‌است که شرایط کشی--ریمان (شرایط مقابل) برقرار باشند:
\EMdisp{\begin{cases}\dfrac{\partial u}{\partial x}=\dfrac{\partial v}{\partial y}\\\noalign{\vskip 2mm} \dfrac{\partial u}{\partial y}=-\dfrac{\partial v}{\partial x}\end{cases}\qquad\EMbr  w\EMbr =u+iv\EMbr =f(x+iy)\EMbr =f(z)}

\end{EMtheorem}

\begin{EMproof}

\EMdisp{f'(z)\EMbr =\lim_{|\Delta z|\to0}\frac{f(z+\Delta z)-f(z)}{\Delta z},\qquad\EMbr  z\EMbr =x+iy,\qquad\EMbr  \Delta z\EMbr =\Delta x+i\Delta y}
\EMdisp{\begin{aligned}
f(z+\Delta z)-f(z)&=\big[u(x+\Delta x,y+\Delta y)+iv(x+\Delta x,y+\Delta y)\big]-\big[u(x,y)+iv(x,y)\big]\\
&=\big[u(x+\Delta x,y+\Delta y)-u(x,y)\big]+i\big[v(x+\Delta x,y+\Delta y)-v(x,y)\big]
\end{aligned}}
\EMdisp{\frac{f(z+\Delta z)-f(z)}{\Delta z}\EMbr =\frac{u(x+\Delta x,y+\Delta y)-u(x,y)}{\Delta x+i\Delta y}+i\,\frac{v(x+\Delta x,y+\Delta y)-v(x,y)}{\Delta x+i\Delta y}}
شرط لازم برای وجود مشتق \EMim{w=f(z)} آن است که اگر از هر جهت به نقطهٔ \EMim{z} نزدیک شویم (از جمله در راستای افقی یا عمودی) مقدار حد مربوط به \EMim{f'(z)} تغییر نکند.

\textbf{مسیر ۱: راستای قائم.} \EMim{\Delta z=\Delta x+i\Delta y} و \EMim{\Delta x=0}، یعنی \EMim{\Delta z=i\Delta y}:
\EMdisp{\begin{aligned}
f'(z)&=\lim_{|\Delta z|\to0}\frac{f(z+\Delta z)-f(z)}{\Delta z}\\
&=\lim_{\Delta y\to0}\left[\frac{u(x,y+\Delta y)-u(x,y)}{i\Delta y}+i\,\frac{v(x,y+\Delta y)-v(x,y)}{i\Delta y}\right]\\
&=\frac1i\,\frac{\partial u}{\partial y}+\frac{\partial v}{\partial y}=\frac{\partial v}{\partial y}-i\,\frac{\partial u}{\partial y}
\end{aligned}}

\textbf{مسیر ۲: راستای افقی.} \EMim{\Delta z=\Delta x+i\Delta y} و \EMim{\Delta y=0}، یعنی \EMim{\Delta z=\Delta x}:
\EMdisp{\begin{aligned}
f'(z)&=\lim_{|\Delta z|\to0}\frac{f(z+\Delta z)-f(z)}{\Delta z}\\
&=\lim_{\Delta x\to0}\left[\frac{u(x+\Delta x,y)-u(x,y)}{\Delta x}+i\,\frac{v(x+\Delta x,y)-v(x,y)}{\Delta x}\right]\\
&=\frac{\partial u}{\partial x}+i\,\frac{\partial v}{\partial x}
\end{aligned}}

بنابراین
\EMdisp{f'(z)\EMbr =\frac{\partial v}{\partial y}-i\,\frac{\partial u}{\partial y}\EMbr =\frac{\partial u}{\partial x}+i\,\frac{\partial v}{\partial x}}
و با مساوی قرار دادن بخش‌های حقیقی و موهومی:

\end{EMproof}

\EMfigure{m02-cr-paths}{دو مسیر نزدیک شدن به \EMim{z}: راستای قائم (\EMim{\Delta z=i\Delta y}) و
راستای افقی (\EMim{\Delta z=\Delta x})}

\begin{EMimportant}{شرایط کشی--ریمان}

\EMdisp{\frac{\partial u}{\partial x}\EMbr =\frac{\partial v}{\partial y},\qquad\EMbr  \frac{\partial u}{\partial y}\EMbr =-\frac{\partial v}{\partial x}}

\end{EMimportant}

\begin{EMremark}{نتیجه}

\begin{enumerate}
\def\labelenumi{\arabic{enumi}.}
\tightlist
\item
  اگر این شرایط برقرار نباشد مشتق در نقطهٔ \EMim{z} وجود ندارد.
\item
  روش‌های زیر جهت مشتق‌گیری پیشنهاد می‌شود:
  \EMdisp{f'(z)\EMbr =\lim_{|\Delta z|\to0}\frac{f(z+\Delta z)-f(z)}{\Delta z}\EMbr =\frac{\partial v}{\partial y}-i\,\frac{\partial u}{\partial y}\EMbr =\frac{\partial u}{\partial x}+i\,\frac{\partial v}{\partial x}\EMbr =\frac{\partial u}{\partial x}-i\,\frac{\partial u}{\partial y}\EMbr =\frac{\partial v}{\partial y}+i\,\frac{\partial v}{\partial x}}
\end{enumerate}

\end{EMremark}

\begin{EMtheorem}{شرط کافی مشتق‌پذیری}

شرط کافی برای آنکه تابع مختلط \EMim{w=f(z)} مشتق‌پذیر باشد آن‌است که شرایط کشی--ریمان برقرار باشند.

اثبات خارج از حوصلهٔ کلاس است.

\end{EMtheorem}

\begin{EMtheorem}{قضیهٔ کشی--ریمان}

شرط لازم و کافی برای آنکه تابع مختلط \EMim{w=f(z)} مشتق‌پذیر باشد آن‌است که شرایط کشی--ریمان (شرایط مقابل) برقرار باشند.
\EMdisp{\begin{cases}\dfrac{\partial u}{\partial x}=\dfrac{\partial v}{\partial y}\\\noalign{\vskip 2mm} \dfrac{\partial u}{\partial y}=-\dfrac{\partial v}{\partial x}\end{cases}\qquad\EMbr  w\EMbr =u+iv\EMbr =f(x+iy)\EMbr =f(z)}

\end{EMtheorem}

\EMsubsection{نتیجهٔ برقراری شرایط کشی--ریمان (تحلیلی بودن یک تابع)}{نتیجهٔ برقراری شرایط کشی–ریمان (تحلیلی بودن یک تابع)}

از شرایط کشی--ریمان نسبت به \EMim{x} و \EMim{y} مشتق می‌گیریم:

\EMdisp{\left.\begin{aligned}\frac{\partial^2u}{\partial x^2}&=\frac{\partial^2v}{\partial x\partial y}\\[1mm] \frac{\partial^2u}{\partial y^2}&=-\frac{\partial^2v}{\partial x\partial y}\end{aligned}\right\}\ \xrightarrow{\ \EMt{جمع}\ }\ \frac{\partial^2u}{\partial x^2}+\frac{\partial^2u}{\partial y^2}\EMbr =0}

\EMdisp{\left.\begin{aligned}\frac{\partial^2u}{\partial x\partial y}&=\frac{\partial^2v}{\partial y^2}\\[1mm] \frac{\partial^2u}{\partial x\partial y}&=-\frac{\partial^2v}{\partial x^2}\end{aligned}\right\}\ \xrightarrow{\ \EMt{تفریق}\ }\ \frac{\partial^2v}{\partial x^2}+\frac{\partial^2v}{\partial y^2}\EMbr =0}

\begin{EMremark}{نتیجه}

قسمت حقیقی و موهومی یک تابع تحلیلی در معادلهٔ لاپلاس دوبعدی صدق می‌کند.

\end{EMremark}

\EMhold

\EMsection{مشتق‌پذیری توابع مختلط: مثال‌ها}{مشتق‌پذیری توابع مختلط: مثال‌ها}

\begin{EMexample}{}

تابع \EMim{w=z^2} را در نظر بگیرید؛ آیا این تابع تحلیلی است؟

\EMdisp{w\EMbr =f(z)\EMbr =z^2\EMbr =(x+iy)^2\EMbr =(x^2-y^2)+i\,2xy,\qquad\EMbr  u(x,y)\EMbr =x^2-y^2,\qquad\EMbr  v(x,y)\EMbr =2xy}
\EMdisp{\frac{\partial u}{\partial x}\EMbr =2x\EMbr =\frac{\partial v}{\partial y}\ \EMcheck \qquad\EMbr \qquad\EMbr  \frac{\partial u}{\partial y}\EMbr =-2y\EMbr =-\frac{\partial v}{\partial x}\ \EMcheck }
\EMdisp{f'(z)\EMbr =\frac{\partial u}{\partial x}+i\,\frac{\partial v}{\partial x}\EMbr =2x+i\,2y\EMbr =2(x+iy)\EMbr =2z}
\EMdisp{\nabla^2u\EMbr =2-2\EMbr =0,\qquad\EMbr  \nabla^2v\EMbr =0+0\EMbr =0}

\end{EMexample}

\begin{EMexample}{}

تابع \EMim{w=|z|^2} را در نظر بگیرید؛ آیا این تابع تحلیلی است؟

\EMdisp{w\EMbr =f(z)\EMbr =|z|^2\EMbr =(x+iy)(x-iy)\EMbr =x^2+y^2,\qquad\EMbr  u(x,y)\EMbr =x^2+y^2,\qquad\EMbr  v(x,y)\EMbr =0}
\EMdisp{\frac{\partial u}{\partial x}\EMbr =2x,\quad\EMbr  \frac{\partial v}{\partial y}\EMbr =0,\qquad\EMbr  \frac{\partial u}{\partial y}\EMbr =2y,\quad\EMbr  -\frac{\partial v}{\partial x}\EMbr =0}
\EMdisp{\nabla^2u\EMbr =2+2\EMbr =4,\qquad\EMbr  \nabla^2v\EMbr =0+0\EMbr =0}
تحلیلی در هیچ‌جا به جز مبدأ.

\end{EMexample}

\begin{EMexercise}{}

بررسی کنید آیا هر کدام از توابع زیر تحلیلی‌اند. در صورت تحلیلی بودن مشتق آن توابع را به دست آورید.

\begin{enumerate}
\def\labelenumi{\arabic{enumi}.}
\tightlist
\item
  \EMim{f(z)=\bar z}
\item
  \EMim{f(z)=\operatorname{Re}\{z\}}
\item
  \EMim{f(z)=z^3+2z^2+3i}
\end{enumerate}

\end{EMexercise}

\EMhold

\EMsection{توابع مختلط تحلیلی}{توابع مختلط تحلیلی}

\begin{EMexample}{}

\EMim{u=x^2-y^2} قسمت حقیقی یک تابع تحلیلی \EMim{f(z)} است. مطلوب است محاسبهٔ \EMim{f(z)}.

از آنجائی که \EMim{f(z)} تحلیلی است، قسمت‌های حقیقی و موهومی تابع در شرایط کوشی--ریمان صدق می‌کنند و در نتیجه این قسمت‌ها در معادلهٔ لاپلاس دو بعدی نیز باید صدق کنند.

۱) اول چک کنیم که آیا قسمت حقیقی تابع در معادلهٔ لاپلاس ۲ بعدی صدق می‌کند؟
\EMdisp{\nabla^2u\EMbr =\frac{\partial^2}{\partial x^2}(x^2-y^2)+\frac{\partial^2}{\partial y^2}(x^2-y^2)\EMbr =2-2\EMbr =0\ \EMcheck }

۲) از روابط کشی--ریمان استفاده می‌کنیم:
\EMdisp{\frac{\partial u}{\partial x}\EMbr =2x\EMbr =\frac{\partial v}{\partial y}\quad\EMbr \EMbr \Longrightarrow\quad\EMbr  v(x,y)\EMbr =2xy+\varphi(x)}
\EMdisp{-\frac{\partial v}{\partial x}\EMbr =-2y-\varphi'(x)\EMbr =\frac{\partial u}{\partial y}\EMbr =-2y\quad\EMbr \EMbr \Longrightarrow\quad\EMbr \varphi'(x)\EMbr =0\quad\EMbr \EMbr \Longrightarrow\quad\EMbr \varphi(x)\EMbr =c,\ \ c\in\mathbb{R}}
\EMdisp{f(z)\EMbr =(x^2-y^2)+i(2xy+c)\EMbr =z^2+C}

\end{EMexample}

\begin{EMexample}{}

اگر \EMim{u=\cos x\cosh y}،
الف) نشان دهید این تابع در صفحه هارمونیک است (\EMim{\nabla^2u=0}).
ب) تابع \EMim{v(x,y)} را چنان بیابید که \EMim{u+iv} تحلیلی شود.

الف) چک کنیم که آیا قسمت حقیقی تابع در معادلهٔ لاپلاس ۲ بعدی صدق می‌کند؟
\EMdisp{\nabla^2u\EMbr =\frac{\partial^2}{\partial x^2}(\cos x\cosh y)+\frac{\partial^2}{\partial y^2}(\cos x\cosh y)\EMbr =-\cos x\cosh y+\cos x\cosh y\EMbr =0\ \EMcheck }

ب) از روابط کشی--ریمان استفاده می‌کنیم:
\EMdisp{\frac{\partial u}{\partial x}\EMbr =-\sin x\cosh y\EMbr =\frac{\partial v}{\partial y}\quad\EMbr \EMbr \Longrightarrow\quad\EMbr  v(x,y)\EMbr =-\sin x\sinh y+\varphi(x)}
\EMdisp{-\frac{\partial v}{\partial x}\EMbr =\cos x\sinh y-\varphi'(x)\EMbr =\frac{\partial u}{\partial y}\EMbr =\cos x\sinh y\quad\EMbr \EMbr \Longrightarrow\quad\EMbr \varphi'(x)\EMbr =0}
\EMdisp{\varphi(x)\EMbr =c,\ \ c\in\mathbb{R}\quad\EMbr \EMbr \Longrightarrow\quad\EMbr  v(x,y)\EMbr =-\sin x\sinh y+c}

\end{EMexample}

\EMsection{شرایط کشی--ریمان در مختصات قطبی}{شرایط کشی–ریمان در مختصات قطبی}

\EMdisp{\begin{cases}x=r\cos\theta\\ y=r\sin\theta\end{cases}\qquad\EMbr \begin{cases}r=\sqrt{x^2+y^2}\\ \theta=\tan^{-1}\dfrac{y}{x}\end{cases}\qquad\EMbr  r_x\EMbr =\frac{x}{r},\quad\EMbr \theta_x\EMbr =\frac{-y}{r^2},\quad\EMbr  r_y\EMbr =\frac{y}{r},\quad\EMbr \theta_y\EMbr =\frac{x}{r^2}}

\EMdisp{\frac{\partial u}{\partial x}\EMbr =\frac{\partial u}{\partial r}\frac{\partial r}{\partial x}+\frac{\partial u}{\partial\theta}\frac{\partial\theta}{\partial x}\EMbr =\frac xr\frac{\partial u}{\partial r}-\frac{y}{r^2}\frac{\partial u}{\partial\theta}\EMbr =\cos\theta\frac{\partial u}{\partial r}-\frac1r\sin\theta\frac{\partial u}{\partial\theta}}

\EMdisp{\frac{\partial v}{\partial y}\EMbr =\frac{\partial v}{\partial r}\frac{\partial r}{\partial y}+\frac{\partial v}{\partial\theta}\frac{\partial\theta}{\partial y}\EMbr =\frac yr\frac{\partial v}{\partial r}+\frac{x}{r^2}\frac{\partial v}{\partial\theta}\EMbr =\sin\theta\frac{\partial v}{\partial r}+\frac1r\cos\theta\frac{\partial v}{\partial\theta}}

\EMdisp{\frac{\partial u}{\partial x}\EMbr =\frac{\partial v}{\partial y}\;\EMbr \Longrightarrow\;\cos\theta\frac{\partial u}{\partial r}-\frac1r\sin\theta\frac{\partial u}{\partial\theta}\EMbr =\sin\theta\frac{\partial v}{\partial r}+\frac1r\cos\theta\frac{\partial v}{\partial\theta}}

مشابهاً

\EMdisp{\frac{\partial u}{\partial y}\EMbr =-\frac{\partial v}{\partial x}\;\EMbr \Longrightarrow\;\sin\theta\frac{\partial u}{\partial r}+\frac1r\cos\theta\frac{\partial u}{\partial\theta}\EMbr =-\cos\theta\frac{\partial v}{\partial r}+\frac1r\sin\theta\frac{\partial v}{\partial\theta}}

\begin{EMimportant}{شرایط کشی--ریمان در فرم قطبی}

\EMdisp{\frac{\partial u}{\partial r}\EMbr =\frac1r\frac{\partial v}{\partial\theta},\qquad\EMbr  -\frac1r\frac{\partial u}{\partial\theta}\EMbr =\frac{\partial v}{\partial r}}

\end{EMimportant}

\begin{EMexercise}{}

با استفاده از شرایط کشی--ریمان در فرم قطبی ثابت کنید که \EMim{u} و \EMim{v} در معادلهٔ لاپلاس در فرم قطبی صدق می‌کنند، یعنی:
\EMdisp{\nabla^2u\EMbr =u_{rr}+\frac1{r^2}u_{\theta\theta}+\frac1ru_r\EMbr =0,\qquad\EMbr  \nabla^2v\EMbr =v_{rr}+\frac1{r^2}v_{\theta\theta}+\frac1rv_r\EMbr =0}

\end{EMexercise}
